\documentclass[11pt]{article}
\usepackage[margin=1.8cm, top=1.5cm, bottom=1.5cm]{geometry}
\usepackage{amsmath, amssymb, amsthm}
\usepackage{microtype}
\usepackage{parskip}
\usepackage[T1]{fontenc}
\usepackage{lmodern}

\newtheorem*{theorem}{Theorem}

\title{\textbf{Any Polygon Is a Sum of Regular Polygons}\\[4pt]
       \large A note on the Fourier decomposition in $\mathbb{C}^n$}
\author{}
\date{}

\begin{document}
\maketitle
\thispagestyle{empty}
\vspace{-1em}

\subsection*{1.\enspace A polygon as a vector}

Label $n$ equally spaced moments in time $0, 1, \ldots, n{-}1$ and record a
complex number $z_j \in \mathbb{C}$ at each one.
The list
\[
  \mathbf{z} = (z_0,\, z_1,\, \ldots,\, z_{n-1}) \;\in\; \mathbb{C}^n
\]
is both a vector in an $n$-dimensional complex space \emph{and} a polygon:
connect the dots $z_0 \to z_1 \to \cdots \to z_{n-1} \to z_0$ in the plane.
Every choice of $n$ vertices gives a different polygon, so polygons and
vectors in $\mathbb{C}^n$ are the same thing.

\subsection*{2.\enspace The regular polygons — a special basis}

Fix a primitive $n$-th root of unity $\omega = e^{2\pi i/n}$.
For each frequency $k = 0, 1, \ldots, n{-}1$ define the vector
\[
  \mathbf{e}_k
  \;=\;
  \bigl(1,\; \omega^k,\; \omega^{2k},\; \ldots,\; \omega^{(n-1)k}\bigr).
\]
The $j$-th vertex of $\mathbf{e}_k$ is $\omega^{jk}$, a point on the unit
circle at angle $\tfrac{2\pi jk}{n}$.
Consecutive vertices are separated by the same angle, so \emph{$\mathbf{e}_k$
is a regular polygon} (or a star polygon when $k > 1$, tracing the $n$
vertices of a regular $n$-gon in steps of $k$).

The $n$ vectors $\mathbf{e}_0, \mathbf{e}_1, \ldots, \mathbf{e}_{n-1}$ are
mutually orthogonal under the standard inner product on $\mathbb{C}^n$:
\[
  \langle \mathbf{e}_j,\, \mathbf{e}_k \rangle
  = \sum_{m=0}^{n-1} \omega^{mj}\,\overline{\omega^{mk}}
  = \sum_{m=0}^{n-1} \omega^{m(j-k)}
  = \begin{cases} n & j = k, \\ 0 & j \neq k. \end{cases}
\]
They therefore form an \textbf{orthogonal basis} of $\mathbb{C}^n$ —
the \emph{Fourier basis}.

\subsection*{3.\enspace Every polygon splits uniquely into regular polygons}

Because $\{\mathbf{e}_k\}$ is a basis, any polygon $\mathbf{z} \in \mathbb{C}^n$
expands uniquely as
\[
  \mathbf{z}
  \;=\;
  \sum_{k=0}^{n-1} \hat{z}_k\, \mathbf{e}_k,
  \qquad
  \hat{z}_k
  \;=\;
  \frac{\langle \mathbf{z},\, \mathbf{e}_k \rangle}{n}
  \;=\;
  \frac{1}{n}\sum_{j=0}^{n-1} z_j\, \omega^{-jk}.
\]
The scalar $\hat{z}_k \in \mathbb{C}$ is the \textbf{$k$-th Fourier
coefficient}.  The term $\hat{z}_k\,\mathbf{e}_k$ is a scaled and rotated
copy of the regular polygon $\mathbf{e}_k$: multiplying by $\hat{z}_k$ scales
every vertex by $|\hat{z}_k|$ and rotates the whole polygon by $\arg(\hat{z}_k)$.
The formula above is the \textbf{Discrete Fourier Transform} (DFT).

\begin{theorem}
Every polygon with $n$ vertices can be written in exactly one way as a sum
of $n$ regular polygons (one per frequency $k = 0, \ldots, n{-}1$).
The $k$-th summand is the orthogonal projection of the polygon onto the
frequency-$k$ direction.
\end{theorem}

\subsection*{4.\enspace What the applet shows}

The applet lets you set the number of vertices $n$ and an exponent $k$.
It draws the regular polygon $\mathbf{e}_k$ — the $k$-th basis vector — as
its starting configuration.
The \emph{exp} control changes $k$, cycling through the basis elements:

\begin{itemize}
  \item $k = 1$: the standard regular $n$-gon, vertices at
        $1, \omega, \omega^2, \ldots$ in order.
  \item $k = 2, 3, \ldots$: the same $n$ points visited in larger steps,
        producing star shapes when $\gcd(k, n) = 1$, or a
        smaller repeated polygon when $\gcd(k,n) > 1$.
  \item $k = 0$: all $n$ vertices collapse to $1$ (zero-frequency, the
        ``DC'' term representing the centroid).
\end{itemize}

The physics simulation then perturbs the shape and lets it oscillate,
giving intuition for how a rigid regular polygon resists deformation — and
how any deformed polygon can be understood as a superposition of these
pure-frequency shapes vibrating independently.

\vfill
\noindent\rule{\linewidth}{0.4pt}\\[2pt]
{\small
$\mathbb{C}^n$ is a complex vector space of dimension $n$.
$\omega = e^{2\pi i/n}$ satisfies $\omega^n = 1$ and $\omega^j \neq 1$ for
$0 < j < n$.
The DFT matrix $F$ with $F_{jk} = \omega^{jk}$ is unitary up to the factor
$1/\sqrt{n}$; its columns are exactly $\mathbf{e}_0, \ldots, \mathbf{e}_{n-1}$.}

\end{document}
